% ===========================================================================
%  utente/04-streams/00-capitolo.tex
% ===========================================================================
\chapter{Streams}\label{cap:streams}

Uno \voce{stream} è il risultato di una strategia: la serie dei suoi
guadagni e delle sue perdite nel tempo. La pagina \voce{Streams} li mostra
tutti, una scheda per ciascuno.

\figuraSchermata{1.0}{figure/schermate/streams.png}
  {L'elenco dei risultati. Ogni scheda porta l'andamento, lo Sharpe, il
   rendimento e la perdita massima.}
  {streams}

\section{Lanciare un backtest}\label{sez:run}

La scheda \voce{Run backtest}, accanto a \voce{Streams}, elenca tutto ciò
che si può far girare: le strategie e gli approcci installati, ciascuno con
la propria descrizione e un pulsante.

\figuraSchermata{1.0}{figure/schermate/run.png}
  {Che cosa si può lanciare. \pulsante{Run All} li mette in coda tutti, uno
   dopo l'altro.}
  {run}

\begin{pratica}{Lanciare, e poi}
\begin{enumerate}
  \item Premere il pulsante della strategia voluta.
  \item La pagina mostra il registro in diretta: si vede che cosa sta
        facendo, riga per riga.
  \item Si può chiudere la pagina. Il lavoro continua, e la campanella in
        alto a destra avvisa alla fine.
  \item Il risultato compare in \voce{Streams}.
\end{enumerate}
\end{pratica}

\begin{attenzione}{Un lavoro per volta}
I backtest girano in fila, non in parallelo, anche premendo più pulsanti.
Non è un limite da aggirare: due calcoli contemporanei si contenderebbero
gli stessi file e produrrebbero risultati mescolati.
\end{attenzione}

\section{La scheda di un risultato}\label{sez:stream-dettaglio}

Cliccando su uno stream si apre la sua pagina.

\figuraSchermata{1.0}{figure/schermate/stream.png}
  {La scheda di uno stream. In alto il giudizio di robustezza, sotto le
   metriche, poi la curva.}
  {stream}

\subsection{Robustness: il riquadro da leggere per primo}

È la parte più importante della pagina, ed è quella che risponde alla
domanda «questa strategia ha funzionato, o ha avuto fortuna?». Il punteggio
a destra riassume i cinque numeri in uno.

\begin{center}
\begin{tabular}{@{}L{42mm}L{96mm}@{}}
\toprule
\textbf{Indicatore} & \textbf{Che cosa dice} \\
\midrule
\rigaTenue
\voce{P(Sharpe > 0)} & La probabilità che il risultato non sia rumore. Sotto
  il \num{95}\,\% la strategia non ha dimostrato niente. \\
\voce{Monte Carlo p-value} & La stessa domanda posta in un altro modo:
  rimescolando i dati, quante volte un caso fortunato avrebbe fatto
  altrettanto bene. Più è piccolo, meglio è. \\
\rigaTenue
\voce{Min. track record} & Da quanto tempo questa strategia dovrebbe girare
  perché il suo risultato sia credibile. Un valore superiore al periodo
  provato significa: non c'è ancora abbastanza storia. \\
\voce{OOS Sharpe instability} & Quanto lo Sharpe cambia da un periodo
  all'altro. Alto significa che la strategia ha funzionato in alcuni periodi
  e non in altri. \\
\rigaTenue
\voce{Losing OOS periods} & La percentuale di periodi chiusi in perdita. \\
\bottomrule
\end{tabular}
\end{center}

\begin{perche}{Perché un giudizio «Doubtful» non è un errore}
Significa che i dati non bastano a distinguere questa strategia dal caso.
Non che sia sbagliata: che non è stata dimostrata. È la risposta onesta su
un periodo breve, e il modo per cambiarla è avere più storia --- non
cambiare i parametri finché il giudizio migliora, che è il modo più comune
di ingannarsi.
\end{perche}

\subsection{Le metriche}

Dodici riquadri. Oltre a quelli già visti nella schermata iniziale
(\vediCap{overview}):

\begin{description}[leftmargin = 1.4em, style = nextline, itemsep = 4pt]
  \item[\voce{Beta vs SPY}] Quanto la strategia segue il mercato. Vicino a
    zero: fa una cosa sua. Vicino a uno: si muove come l'indice.
  \item[\voce{Sortino}] Come lo Sharpe, ma conta solo le oscillazioni verso
    il basso. Salire in fretta non è un difetto.
  \item[\voce{Calmar}] Rendimento annuo diviso perdita massima. Risponde a
    «quanto rende per ogni punto di sofferenza».
  \item[\voce{Hit rate}] La percentuale di periodi in guadagno.
  \item[\voce{N obs}] Quante osservazioni. Poche osservazioni rendono tutto
    il resto meno affidabile.
\end{description}

\subsection{I tre pulsanti in alto}

\begin{itemize}
  \item \pulsante{Open the code in Studio} --- apre l'artefatto nell'editor.
  \item \pulsantesec{Run backtest} --- rilancia il calcolo.
  \item \pulsantesec{Use in Ensemble} --- porta questo stream nel
    costruttore dei project (\vediCap{projects}).
\end{itemize}

\section{View: guardare che cosa ha fatto}\label{sez:view}

La voce \voce{View} del menu mostra i prezzi con sopra le operazioni della
strategia: dove è entrata, dove è uscita, e gli indicatori che stava
guardando.

\figuraSchermata{1.0}{figure/schermate/visual.png}
  {La vista sui grafici. Si sceglie la strategia e lo strumento, e si
   sovrappongono gli indicatori.}
  {visual}

\noindent
Serve a una cosa che i numeri non dicono: \emph{se la strategia sta facendo
quello che si pensava}. Una curva che sale bene può nascondere venti
operazioni sensate e una sola enorme fortunata.
